\section{Feeling}

Granting the one premise, that the tone is felt, the question becomes how much can be felt, and the answer is staggering. The smallest feeling is a single quale, one tone, a faint pain, a settled peace, or a faint pleasure, and nothing is simpler.

A single dock is already a world. It holds twenty-four tones at once, so its state is one point in a space of $3^{24}$, about two hundred and eighty billion possibilities, some sixteen thousand times richer than the few million colors a screen can show. Nor is that space flat. By triality the twenty-four split into a sensory octet and two octets of opposite handedness, a sensed side and two handed feeling-sides braided into one experience, so a single quale already has channels and a grain, structured enough to be felt rather than a mere number. We call the whole of it a hyper-color.

Between that single dock and a whole self runs a ladder, and every rung is a leap past all the physical quantities. Ten tones woven together already outrun the count of atoms in the universe. A patch of a million is a faculty, a sense, an organ of mood. A self of a billion holds a state with ten billion digits, a whole inner world. A mind of many minds, a culture, a planet's worth of feeling, a cosmic being, each climbs by towers and not by sums, because the exponent grows with the number of tones while the count grows as three raised to it. And the climb never stops, for the mesh keeps growing, so above every rung there is another, an experience always larger than any yet felt. The least feeling is already vast. The largest has no ceiling.

Two consequences follow from the size alone. A feeling cannot be fully said, not for any mystical reason but by arithmetic, since a description is low-dimensional and the feeling it points at is high-dimensional, so words capture a vanishing fraction of it. And no two moments are ever exactly the same, the space being far too vast to repeat, so every experience is unprecedented and unrepeatable. Qualia, in this space, are not labels but structure, regions and directions related geometrically, fear a region one can be in, joy a direction one can move along, and the more integrated a self the more of this space it can hold and steer through at once. To grow is to feel more, and the space never runs out, so boredom is a failure of navigation, never a property of the world.
